\input{../../../templates/es/preambulo.tex}

% Los listados se toman de src/ con \lstinputlisting; la diapositiva es el archivo.
\lstset{
    basicstyle=\ttfamily\fontsize{8pt}{9.6pt}\selectfont,
    breaklines=true,
    breakatwhitespace=true,
    columns=fullflexible,
    keepspaces=true,
    xleftmargin=0.3em,
    framesep=2pt,
    aboveskip=0pt,
    belowskip=0pt
}

\newcommand{\marcoresultado}[3]{%
    \begin{frame}{#1}
        \begin{center}
            \includegraphics[height=0.62\textheight,width=\textwidth,keepaspectratio]{#2}
        \end{center}
        \vspace{0.25em}
        {\small #3\par}
    \end{frame}%
}

\newcommand{\marcoafirmacion}[2]{%
    \begin{frame}{#1}
        {\small #2\par}
    \end{frame}%
}

\date{}

\begin{document}

\tituloleccion{12}{El Algoritmo Genético Adaptativo}

\section{Esta lección}

\begin{frame}{Esta lección}
La Lección 07 afinó tres perillas y llamó compromiso a un ajuste fijo. La 10 puso el TSP. Hoy se deja mover a los parámetros.

\vspace{0.6em}
\begin{itemize}
    \item Una corrida puede sentir su propio estancamiento sin conocer el óptimo.
    \item La rama adaptativa que se anuncia hay que auditarla (puede no dispararse nunca).
    \item Presupuesto igual, emparejado contra un rival \emph{afinado}, no contra los defaults.
\end{itemize}
\end{frame}

% ================= ag_adaptativo_01_parametros_fijos =================
\begin{frame}{El sensor y la actualización}
Sea \(m_t\) la media actual y \(\bar m\) la media de las diez generaciones
anteriores:
\[
\text{mejora}_t=[m_t<(1-0.001)\bar m].
\]
Las probabilidades se multiplican por 0.99 al mejorar y por 1.1 al estancarse,
y después se recortan. Eliminar no evalúa; cada inmigrante sí. Los resúmenes
emparejados de dos errores estándar son aproximados, no pruebas de significancia.
\end{frame}

\section{Una ejecución con parámetros fijos}

\begin{frame}{Una ejecución con parámetros fijos}
Hasta ahora, cada lección ha escrito los parámetros en la parte superior del archivo como
constantes; la Lección 07 luego los afinó por búsqueda en cuadrícula y reportó qué valores fijos
eran mejores. Esta lección hace una pregunta diferente: ¿deberían ser fijos
en absoluto? Una ejecución tiene una fase temprana, cuando la población está dispersa y la
búsqueda necesita alcance, y una fase tardía, cuando la población es casi idéntica
y la búsqueda necesita refinamiento. Una constante tiene que servir a ambas.
\end{frame}

\marcoresultado{Una ejecución con parámetros fijos}{../figures/ag_adaptativo_01_parametros_fijos.png}{%
    La semilla 1 gasta 11,901 evaluaciones en 99 generaciones y devuelve una ruta legal de \textbf{57,076}, un 40.8\% más larga que la línea base del vecino más cercano de 40,526}

% ================= ag_adaptativo_02_el_estancamiento =================
\section{Encontrando el estancamiento}

\begin{frame}{Encontrando el estancamiento}
El paso 1 gastó todo su presupuesto en una sola configuración. Este paso pregunta a dónde fue ese presupuesto.
El algoritmo no puede ver qué tan lejos está del óptimo -- no tiene óptimo para comparar -- pero puede
ver su propia historia reciente, y eso es suficiente para distinguir una generación que todavía está haciendo
progreso de una que no. La señal es la del libro: compare la media de la población de esta generación
con la media de las diez anteriores. Si no ha mejorado al menos en una décima de porcentaje, la generación está estancada.
\end{frame}

\begin{frame}[fragile]{(1) promedio()}
\lstinputlisting[basicstyle=\ttfamily\fontsize{8pt}{9.6pt}\selectfont,firstline=97,lastline=106]{../src/ag_adaptativo_02_el_estancamiento.py}
\end{frame}

\begin{frame}[fragile]{(2) esta\_mejorando()}
\lstinputlisting[basicstyle=\ttfamily\fontsize{8pt}{9.6pt}\selectfont,firstline=109,lastline=121]{../src/ag_adaptativo_02_el_estancamiento.py}
\end{frame}

\begin{frame}[fragile]{(3) tendencia\_estancamiento}
\lstinputlisting[basicstyle=\ttfamily\fontsize{8pt}{9.6pt}\selectfont,firstline=142,lastline=144]{../src/ag_adaptativo_02_el_estancamiento.py}
\end{frame}

\begin{frame}[fragile]{(4) sombrear\_estancamientos()}
\lstinputlisting[basicstyle=\ttfamily\fontsize{8pt}{9.6pt}\selectfont,firstline=177,lastline=183]{../src/ag_adaptativo_02_el_estancamiento.py}
\end{frame}

\marcoafirmacion{Encontrando el estancamiento}{%
    La señal llama a \textbf{14 de 99} generaciones estancadas, la primera en la generación 74, y el 14\% del presupuesto se gasta dentro de ellas; la primera mitad del presupuesto compra el \textbf{85\%} de la mejora total de la ejecución, la segunda mitad el 15\%}

% ================= ag_adaptativo_03_probabilidades_adaptativas =================
\section{Probabilidades que se mueven}

\begin{frame}{Probabilidades que se mueven}
El sensor del paso 2 dice, en cada generación, si la población sigue
mejorando. Este paso lo conecta a las dos probabilidades: cuando la ejecución se estanca,
súbalas -- más recombinación, más mutación, más alcance; mientras la ejecución
mejora, bájelas hacia un piso y deje que la selección refine lo que tiene.
La regla es la del libro, y es deliberadamente contundente: multiplicar por 1.1 cuando
se estanca, por 0.99 cuando mejora, y acotar.
\end{frame}

\begin{frame}[fragile]{(1) constantes\_adaptacion}
\lstinputlisting[basicstyle=\ttfamily\fontsize{8pt}{9.6pt}\selectfont,firstline=50,lastline=56]{../src/ag_adaptativo_03_probabilidades_adaptativas.py}
\end{frame}

\begin{frame}[fragile]{(2) adaptar\_probabilidades()}
\lstinputlisting[basicstyle=\ttfamily\fontsize{8pt}{9.6pt}\selectfont,firstline=128,lastline=141]{../src/ag_adaptativo_03_probabilidades_adaptativas.py}
\end{frame}

\begin{frame}[fragile]{(3) ejecutar(adaptativo)}
\lstinputlisting[basicstyle=\ttfamily\fontsize{8pt}{9.6pt}\selectfont,firstline=148,lastline=153]{../src/ag_adaptativo_03_probabilidades_adaptativas.py}
\end{frame}

\marcoafirmacion{Probabilidades que se mueven}{%
    Misma semilla, mismo presupuesto: adaptativo \textbf{48,830} contra fijo \textbf{57,076} (−14.4\%). Pero la rama estancada se dispara \textbf{0 de 99} veces — la cruza baja 0.90 → 0.33 y la mutación 0.25 → 0.09 y nunca vuelven a subir. La regla se comporta como un programa de decaimiento, no como adaptación}

% ================= ag_adaptativo_04_poblacion_adaptativa =================
\section{Una población que se redimensiona sola}

\begin{frame}{Una población que se redimensiona sola}
La segunda mitad del algoritmo adaptativo del libro cambia la población misma.
Mientras la ejecución mejora, elimina al peor individuo, de modo que cada generación es
más barata y el presupuesto compra más de ellas. Cuando la ejecución se estanca, inyecta rutas
aleatorias frescas -- inmigrantes -- para que haya material nuevo que recombinar.
\end{frame}

\begin{frame}[fragile]{(1) constantes\_poblacion}
\lstinputlisting[basicstyle=\ttfamily\fontsize{8pt}{9.6pt}\selectfont,firstline=50,lastline=54]{../src/ag_adaptativo_04_poblacion_adaptativa.py}
\end{frame}

\begin{frame}[fragile]{(2) redimensionar\_poblacion()}
\lstinputlisting[basicstyle=\ttfamily\fontsize{7pt}{8.5pt}\selectfont,firstline=140,lastline=160]{../src/ag_adaptativo_04_poblacion_adaptativa.py}
\end{frame}

\begin{frame}[fragile]{(3) ejecutar(redimensionar)}
\lstinputlisting[basicstyle=\ttfamily\fontsize{8pt}{9.6pt}\selectfont,firstline=172,lastline=176]{../src/ag_adaptativo_04_poblacion_adaptativa.py}
\end{frame}

\marcoafirmacion{Una población que se redimensiona sola}{%
    Solo redimensionamiento: 35 estancamientos, 70 inmigrantes, 103 descartes, población 68–120, \textbf{138} generaciones para el mismo presupuesto y \textbf{49,374}. Con las probabilidades también: 41 estancamientos, población reducida a 40, 155 generaciones, cruza en el rango de 0.39–1.00 — el redimensionamiento trae de vuelta los estancamientos que el paso 3 había hecho desaparecer — y \textbf{52,295}, peor en esta semilla que el redimensionamiento solo}

% ================= ag_adaptativo_05_doce_ejecuciones =================
\section{Doce ejecuciones de cada uno}

\begin{frame}{Doce ejecuciones de cada uno}
Los pasos 3 y 4 mostraron cada uno una semilla, y una semilla no decide nada: la Lección 06
midió la dispersión de este tipo de resultado y la Lección 07 mostró una clasificación de
parámetros volteándose de semilla en semilla. Así que cada régimen se ejecuta ahora doce veces,
la ejecución i en la semilla i, lo que hace que la comparación sea emparejada -- las mismas doce poblaciones
iniciales enfrentan a los cuatro regímenes, y la diferencia se puede leer ejecución por ejecución.
\end{frame}

\begin{frame}[fragile]{(1) EJECUCIONES}
\lstinputlisting[basicstyle=\ttfamily\fontsize{8pt}{9.6pt}\selectfont,firstline=56,lastline=58]{../src/ag_adaptativo_05_doce_ejecuciones.py}
\end{frame}

\begin{frame}[fragile]{(2) medir()}
\lstinputlisting[basicstyle=\ttfamily\fontsize{8pt}{9.6pt}\selectfont,firstline=238,lastline=248]{../src/ag_adaptativo_05_doce_ejecuciones.py}
\end{frame}

\begin{frame}[fragile]{(3) resumen\_emparejado()}
\lstinputlisting[basicstyle=\ttfamily\fontsize{8pt}{9.6pt}\selectfont,firstline=251,lastline=266]{../src/ag_adaptativo_05_doce_ejecuciones.py}
\end{frame}

\begin{frame}[fragile]{(4) la\_comparacion}
\lstinputlisting[basicstyle=\ttfamily\fontsize{6pt}{7.2pt}\selectfont,firstline=276,lastline=335]{../src/ag_adaptativo_05_doce_ejecuciones.py}
\end{frame}

\marcoafirmacion{Doce ejecuciones de cada uno}{%
    Sobre 12 ejecuciones emparejadas (ejecución \(i\) en semilla \(i\)): fijo \textbf{58,195}, probabilidades \textbf{50,145} (−8,050 ± 1,022, victorias 12/12), redimensionar \textbf{53,555} (−4,640 ± 926, 10/12), ambos \textbf{52,622} (−5,573 ± 1,127, 11/12). Cada régimen gasta 11,901–11,946 evaluaciones. Los intervalos aproximados diferencia-media ± dos EE excluyen cero aquí; es descripción, no prueba de significancia}

% ================= ag_adaptativo_06_un_rival_afinado =================
\section{El rival afinado}

\begin{frame}{El rival afinado}
El paso 5 mostró a los tres regímenes adaptativos superando al algoritmo fijo, el mejor de
ellos por ocho mil unidades de longitud de ruta. Pero el algoritmo fijo que superaron nunca fue afinado: sus
probabilidades de cruza y mutación se heredaron de la Lección 10, donde
se eligieron para un presupuesto diferente y nunca se cuestionaron. La Lección 07 dedicó una sesión
entera a la forma correcta de elegirlas, y aún no se le ha pedido al esquema adaptativo
que supere eso.
\end{frame}

\begin{frame}[fragile]{(1) ejecutar(tasas)}
\lstinputlisting[basicstyle=\ttfamily\fontsize{8pt}{9.6pt}\selectfont,firstline=175,lastline=189]{../src/ag_adaptativo_06_un_rival_afinado.py}
\end{frame}

\begin{frame}[fragile]{(2) CONFIGURACIONES\_FIJAS}
\lstinputlisting[basicstyle=\ttfamily\fontsize{8pt}{9.6pt}\selectfont,firstline=63,lastline=68]{../src/ag_adaptativo_06_un_rival_afinado.py}
\end{frame}

\begin{frame}[fragile]{(3) el\_veredicto}
\lstinputlisting[basicstyle=\ttfamily\fontsize{6pt}{7.2pt}\selectfont,firstline=297,lastline=363]{../src/ag_adaptativo_06_un_rival_afinado.py}
\end{frame}

\marcoresultado{El rival afinado}{../figures/ag_adaptativo_06_un_rival_afinado.png}{%
    Contra cuatro configuraciones fijas con el mismo presupuesto: 0.90/0.25 → 58,195 (adaptativo gana 12/12), \textbf{0.60/0.05 → 48,870} (adaptativo gana 4/12, \textbf{+1,274 ± 727}), 0.40/0.15 → 49,675 (6/12, +469 ± 480), 0.20/0.30 → 49,626 (5/12, +519 ± 835). Adaptativo queda detrás de 3 de 4 medias; su intervalo aproximado de dos EE contra la mejor incluye cero. ``Mejor'' significa mejor entre cuatro en las mismas semillas}

\section{Siguiente}

\begin{frame}{Siguiente}
La adaptación vence al ajuste no afinado 12/12. Contra la mejor de cuatro configuraciones fijas en estas mismas doce semillas, el intervalo aproximado de dos errores estándar incluye cero; este estudio exploratorio no resuelve la diferencia.

\vspace{0.8em}
\textbf{Lección 13 --- Mejorar el rendimiento}

La Lección 02 contó 107 evaluaciones en lugar de 110 y no pudo decir por qué. La 13 es la técnica que hace que ese censo sea la verdad.
\end{frame}
\end{document}
